\documentclass[10pt,a4paper]{amsart}
\usepackage[margin=19mm]{geometry}
\usepackage{amsmath,amssymb,mathtools}
\usepackage[scr=rsfs,cal=euler]{mathalfa}
\usepackage{xfrac}
\usepackage[unicode,hidelinks,bookmarks=false]{hyperref}
\newcommand{\ie}{i.e.\ }
\newcommand{\cf}{cf.\ }
\newcommand{\resp}{respectively\ }
\newcommand{\Subsection}{\subsection}
\providecommand{\nobreakdash}{\nobreak\hbox{-}}
\setlength{\emergencystretch}{2em}
\multlinegap0pt
\newcommand{\SL}{\operatorname{SL}}
\usepackage{amssymb}
\usepackage{mathtools}
\usepackage[shortlabels]{enumitem}
\usepackage{xcolor}
\newtheorem{lemma}{Lemma}
\newcommand{\F}{\mathbb{F}}
\newcommand{\PGL}{\operatorname{PGL}}
\newcommand{\PP}{\mathbb P}
\newcommand{\Q}{\mathbb{Q}}
\newcommand{\Z}{\mathbb{Z}}
\title{The groups $\SL_2(\F_{13})$ and $\SL_2(\F_{19})$ are Galois over $\Q$}
\author{Emir Eray Karabiyik}
\date{Source: arXiv:2609.03946v1 (CC BY 4.0)}
\begin{document}
\maketitle

\noindent\emph{Source and licence.} The groups $\SL_2(\F_{13})$ and $\SL_2(\F_{19})$ are Galois over $\Q$. Version arXiv:2609.03946v1 (CC BY 4.0), distributed by arXiv under the Creative Commons Attribution 4.0 International licence, http://creativecommons.org/licenses/by/4.0/, as recorded in the arXiv licence metadata for this exact version and shown on the abstract page https://arxiv.org/abs/2609.03946v1. Full licence notice: \texttt{licenses/arxiv-2609.03946-CC-BY-4.0.txt}.\par\smallskip

\noindent\textit{Editorial scope.} Counted unit: the article's lemma lem:l19-hecke-cosets (source0.tex lines 893--903). The article's complete original proof of it is delivered with it (source0.tex lines 905--949, one \texttt{proof} environment of 45 lines). No mathematics is written, repaired or re-derived: the statement and the proof are the article's own lines, in the article's own notation, and no retained line is substituted or re-flowed.  Retained uncounted as the source-local context the proof needs: lines 884--886.  Nothing was omitted from any retained span, so the package carries no omission record.  No retained line contains a citation, so the wrapper carries no bibliography and no bibliography entry.  No retained line contains a figure environment, an image-inclusion command or drawing code, so no image is delivered and the package declares no asset dependency.  Editorial formatting only, in addition: an AMS \texttt{amsart} preamble that restates the article's own declarations and macros used by the retained lines, namely the environments \texttt{lemma} on separate counters, together with the standard packages those lines need.  The retained environments print in this document as Lemma 1 in order of appearance, whereas the article prints the same items with the numbering of its own full document.  The complete native source of the article as distributed in the arXiv e-print for this exact version is bundled unchanged as \texttt{sources/209328/source0.tex}.
\begin{equation}\label{eq:yang-defining-product}
 P_0(T,z)=\prod_{c\in\mathcal C_{19}}(z-z_c)\in\Q(T)[z].
\end{equation}

\begin{lemma}\label{lem:l19-hecke-cosets}
Choose a splitting $\mathcal O\otimes\Z_{19}\simeq M_2(\Z_{19})$. It
gives natural equivariant bijections
\[
 \mathcal C_{19}\ \simeq\
 \PGL_2(\F_{19})/B(\F_{19})\ \simeq\ \PP^1(\F_{19}).
\] 
Under these bijections, analytic continuation along a loop around an elliptic
point permutes the roots $z_c$ through the image in $\PGL_2(\F_{19})$ of a
generator of the corresponding elliptic stabilizer.
\end{lemma}

\begin{proof}

Put $\mathcal O_{19}=\mathcal O\otimes\Z_{19}$, and let
  $\mathcal E_{19}\subset\mathcal O_{19}$ be the local Eichler order of level
  $19$. Define
  \[
   K=\mathcal O_{19}^{\times}/\Z_{19}^{\times},
   \qquad
   K_0(19)=\mathcal E_{19}^{\times}/\Z_{19}^{\times}.
  \]  
Reduction modulo $19$ maps $K$ onto
$\PGL_2(\F_{19})$, and $K_0(19)$ is the inverse image of the upper triangular
Borel subgroup $B(\F_{19})$.  The kernel of reduction is contained in
$K_0(19)$, so reduction induces a bijection
\[
 K/K_0(19)\ \simeq\
 \PGL_2(\F_{19})/B(\F_{19}).
\]
The cosets $\mathcal C_{19}$ indexing Yang's double-coset product are right
cosets, and hence form the set $K_0(19)\backslash K$.  We identify this set
with the left-coset space above by inversion:
\[
 K_0(19)g\longmapsto g^{-1}K_0(19).
\]
Under this identification, right multiplication by $k$ corresponds to left
multiplication by $k^{-1}$.  We use this left-coset convention below.
Let $L_0$ be the line spanned by $(1,0)$ in $\F_{19}^2$.  Its stabilizer in
$\PGL_2(\F_{19})$ is $B(\F_{19})$, so
\[
 gB(\F_{19})\longmapsto gL_0
\]
is a well-defined equivariant bijection from the displayed coset space to
the set of lines in $\F_{19}^2$, namely $\PP^1(\F_{19})$.  In particular,
the three sets in the lemma each have twenty elements.

Formula~\eqref{eq:yang-defining-product} attaches the factor $z-z_c$ to the
coset $c$.  A small positively oriented loop around an elliptic point acts on
the Hecke correspondence by permuting its cosets.  This action is induced by
a generator of the corresponding elliptic stabilizer, so analytic
continuation carries the factor $z-z_c$ to the factor indexed by the image of
$c$.
Under the chosen splitting and the inversion convention above, reduction
modulo $19$ turns this coset action into the projective action stated in the
lemma.
\end{proof}

\end{document}
